\chapter{Hot Weather}
\index{Hot Weather}

\section*{Definition and Overview}
Hot weather presents significant health risks, particularly during heatwaves when temperatures remain high for several days. The body normally regulates its temperature by sweating and increasing blood flow to the skin, but extreme heat can overwhelm these systems, leading to heat exhaustion, heat stroke, and worsening of chronic illness. Certain groups are especially vulnerable: older adults, young children, people with chronic disease, outdoor workers, and those without access to cooling. Most heat-related illness is preventable with simple precautions, and knowing the warning signs allows early treatment.

\section*{Epidemiology}
Heatwaves cause more deaths in many countries than any other natural hazard, and hospital admissions rise sharply during hot spells. Older adults account for most heat-related deaths, and the risk rises in people with heart, lung, kidney, and mental health conditions. Outdoor workers, athletes, and people without air-conditioning are at higher risk. With climate change, hot days and heatwaves are becoming more frequent and intense, increasing the importance of heat awareness and preparedness across the community.

\section*{Aetiology and Pathophysiology}
\begin{itemize}
    \item \textbf{Heat gain:} high ambient temperature, direct sun, and physical exertion increase body heat
    \item \textbf{Heat loss:} sweating and skin blood flow remove heat, but humidity, dehydration, and poor airflow reduce their effectiveness
    \item \textbf{Dehydration:} fluid loss through sweating reduces blood volume and cooling capacity
    \item \textbf{Vulnerable physiology:} older adults have reduced thirst, sweating, and cardiovascular reserve; infants cannot regulate temperature efficiently
    \item \textbf{Medications:} diuretics, blood pressure drugs, and some psychiatric medicines impair heat tolerance
    \item \textbf{Pathophysiology:} rising core temperature strains the heart and kidneys and, above 40 degrees Celsius, damages organs directly
    \item \textbf{Heat illness spectrum:} cramps, oedema, syncope, exhaustion, and stroke
\end{itemize}

\section*{Clinical Features}
\begin{itemize}
    \item Heat cramps: painful muscle spasms after exertion in heat
    \item Heat oedema: swelling of the hands and feet
    \item Heat syncope: fainting on standing after heat exposure
    \item Heat exhaustion: heavy sweating, weakness, dizziness, headache, nausea, and rapid pulse
    \item Heat stroke: very high temperature, confusion, seizures, hot dry skin, and collapse, an emergency
    \item Worsening of heart, lung, and kidney conditions
    \item Disturbed sleep and irritability
\end{itemize}

\section*{Differential Diagnosis}
\begin{itemize}
    \item Fever from infection is distinguished from hyperthermia by history and examination
    \item Cardiac events, stroke, and hypoglycaemia can mimic heat collapse
    \item Gastroenteritis and dehydration share symptoms
    \item Medication effects may cause dizziness in the heat
    \item Heat exposure context and temperature measurement guide the diagnosis
\end{itemize}

\section*{When to Seek Medical Care}
Seek medical advice for heat illness that does not improve with cooling and fluids, and for symptoms in vulnerable people. Any suspicion of heat stroke is an emergency.

\textbf{Contact your local emergency services for:}
\begin{itemize}
    \item Confusion, agitation, seizures, or unconsciousness in the heat
    \item Body temperature of 40 degrees Celsius or higher
    \item Hot, dry skin with inability to sweat
    \item Collapse, rapid breathing, or chest pain
    \item A baby, older person, or person with chronic illness who is severely unwell
    \item Failure to improve after 30 minutes of cooling and fluids
\end{itemize}

\section*{Diagnosis and Investigations}
\begin{itemize}
    \item \textbf{Clinical assessment:} heat exposure history and symptoms
    \item \textbf{Temperature measurement:} core temperature for suspected heat stroke
    \item \textbf{Blood tests:} electrolytes and kidney function in moderate to severe illness
    \item \textbf{Hydration assessment:} clinical evaluation
    \item \textbf{Monitoring:} for organ dysfunction in hospitalised patients
    \item \textbf{Heatwave warnings:} Bureau of Meteorology heat warnings guide community response
\end{itemize}

\section*{Management}
\begin{itemize}
    \item \textbf{First aid:} move to shade or air-conditioning, remove excess clothing, and cool with fans, water, and compresses
    \item \textbf{Rehydration:} cool water or oral rehydration solution
    \item \textbf{Heat cramps:} rest, cool down, rehydrate, and stretch gently
    \item \textbf{Heat stroke:} immediate aggressive cooling and emergency hospital care
    \item \textbf{Chronic disease care:} adjust medication under medical advice and monitor symptoms
    \item \textbf{Community response:} cool refuges, check on vulnerable people, and heatwave plans
    \item \textbf{Home cooling:} fans, shut curtains, and open windows at night
\end{itemize}

\section*{Complications}
\begin{itemize}
    \item Heat stroke with brain, kidney, and liver damage
    \item Heart attacks, strokes, and heart failure during hot weather
    \item Kidney injury from dehydration
    \item Worsening of respiratory and diabetic conditions
    \item Falls and injuries in older people
    \item Death, particularly among older adults during heatwaves
\end{itemize}

\section*{Prognosis and Outcomes}
Heat illness responds rapidly to early cooling and rehydration, and most people recover fully. Heat stroke is life-threatening but has a good outcome when cooling begins immediately. Heat-related deaths are almost entirely preventable with preparation: acting on warnings, staying cool, hydrating, and checking on vulnerable people. With sensible precautions, hot weather can be managed safely by everyone, including those at highest risk.

\section*{Prevention and Self-Care}
\begin{itemize}
    \item Follow heatwave warnings and heat health alerts
    \item Stay in air-conditioned or cooled places during extreme heat
    \item Drink water regularly without waiting for thirst
    \item Wear light, loose clothing and a hat, and use sunscreen
    \item Avoid strenuous activity during the hottest hours
    \item Never leave children, older people, or pets in parked cars
    \item Check on older relatives, neighbours, and vulnerable people daily
    \item Learn heat illness first aid and act early
\end{itemize}

\section*{Sources and Further Reading}
The following reputable sources provide further information for healthcare students and patients seeking reliable, current advice about hot weather.

\begin{itemize}
    \item Healthdirect. \textit{Staying safe in hot weather.} https://www.healthdirect.gov.au/hot-weather
    \item Better Health Channel. \textit{Beat the heat in hot weather.} https://www.betterhealth.vic.gov.au/
    \item Australian Government Bureau of Meteorology. \textit{Heatwave warnings.} http://www.bom.gov.au/
    \item World Health Organization. \textit{Heat and health.} https://www.who.int/
    \item NHS. \textit{Heatwave: how to cope in hot weather.} https://www.nhs.uk/live-well/
    \item Centers for Disease Control and Prevention. \textit{Extreme heat.} https://www.cdc.gov/
\end{itemize}
